\documentclass{article}
\usepackage{amsmath,amssymb,amsthm}
\usepackage{algorithm}
\usepackage{algpseudocode}
\usepackage{enumitem}
\usepackage{geometry}
\geometry{margin=1in}
\newtheorem{theorem}{Theorem}[section]
\newtheorem{lemma}{Lemma}[section]
\newtheorem{assumption}{Assumption}[section]
\newtheorem{definition}{Definition}[section]
\newcommand{\R}{\mathbb{R}}
\newcommand{\cC}{\mathcal{C}}
\newcommand{\gap}{\operatorname{gap}}
\newcommand{\dist}{\operatorname{dist}}
\newcommand{\diam}{\operatorname{diam}}
\newcommand{\lin}{\operatorname{lin}}
\newcommand{\conv}{\operatorname{conv}}
\newcommand{\argmin}{\operatorname*{arg\,min}}
\newcommand{\argmax}{\operatorname*{arg\,max}}
\newcommand{\inner}[2]{\langle #1,#2\rangle}
\begin{document}

\section*{Away‑Step Frank‑Wolfe (AFW) Algorithm}

\section*{Mathematical Formulation}
Let $f:\R^{d}\rightarrow\R$ be convex and $L$‑smooth on a compact convex polytope 
\[
\cC=\operatorname{conv}\{v^{1},\dots ,v^{m}\}\subset\R^{d}.
\]
Denote by $\nabla f(x)$ the gradient at $x$.
Given the current iterate $x_{k}\in\cC$ we maintain an \emph{active set}
\[
\cS_{k}\subseteq\{v^{1},\dots ,v^{m}\},\qquad 
x_{k}= \sum_{v\in \cS_{k}} \lambda_{k}^{v} v,\;\; \lambda_{k}^{v}\ge 0,\;\; \sum_{v\in \cS_{k}}\lambda_{k}^{v}=1 .
\]

\begin{enumerate}[label=\arabic*.]
\item \textbf{Frank–Wolfe (FW) direction:}
\[
s_{k}= \argmin_{s\in\cC}\inner{\nabla f(x_{k})}{s},\qquad d_{k}^{\text{FW}}:= s_{k}-x_{k}.
\]
\item \textbf{Away direction:}
\[
v_{k}= \argmax_{v\in \cS_{k}}\inner{\nabla f(x_{k})}{v},\qquad 
d_{k}^{\text{A}}:= x_{k}-v_{k}.
\]
\item \textbf{Choose descent direction:}
\[
d_{k}= 
\begin{cases}
d_{k}^{\text{FW}}, & \text{if }\inner{\nabla f(x_{k})}{d_{k}^{\text{FW}}}\le 
\inner{\nabla f(x_{k})}{d_{k}^{\text{A}}},\\[4pt]
d_{k}^{\text{A}}, & \text{otherwise.}
\end{cases}
\]
\item \textbf{Step‑size:}
\[
\gamma_{\max }=
\begin{cases}
1, & d_{k}=d_{k}^{\text{FW}},\\[4pt]
\displaystyle\frac{\lambda_{k}^{v_{k}}}{1-\lambda_{k}^{v_{k}}}, & d_{k}=d_{k}^{\text{A}},
\end{cases}
\qquad 
\gamma_{k}= \argmin_{\gamma\in[0,\gamma_{\max}]}\; f\bigl(x_{k}+ \gamma d_{k}\bigr).
\]
(Exact line‑search is used; a backtracking rule with the same bound $\gamma_{\max}$ is also admissible.)
\item \textbf{Update iterate and active set:}
\[
x_{k+1}=x_{k}+ \gamma_{k} d_{k}.
\]
If $d_{k}=d_{k}^{\text{FW}}$ we add $s_{k}$ to $\cS_{k}$ (or increase its coefficient);
if $d_{k}=d_{k}^{\text{A}}$ we decrease the coefficient of $v_{k}$, and if it becomes zero we drop $v_{k}$ from $\cS_{k+1}$.
\end{enumerate}

\section*{Pseudocode}
\begin{algorithm}[H]
\caption{Away‑Step Frank‑Wolfe (AFW)}
\begin{algorithmic}[1]
\State \textbf{Input:} convex $L$‑smooth $f$, polytope $\cC=\conv\{v^{1},\dots ,v^{m}\}$, tolerance $\varepsilon>0$.
\State Choose an extreme point $v^{i_{0}}\in\cC$, set $x_{0}=v^{i_{0}}$, $\cS_{0}=\{v^{i_{0}}\}$, $\lambda^{v^{i_{0}}}_{0}=1$.
\For{$k=0,1,2,\dots$}
    \State Compute gradient $g_{k}= \nabla f(x_{k})$.
    \State $s_{k}= \argmin_{s\in\cC}\inner{g_{k}}{s}$   \Comment{FW vertex}
    \State $v_{k}= \argmax_{v\in \cS_{k}}\inner{g_{k}}{v}$ \Comment{Away vertex}
    \State $d_{k}^{\text{FW}}=s_{k}-x_{k}$,\; $d_{k}^{\text{A}}=x_{k}-v_{k}$.
    \If{$\inner{g_{k}}{d_{k}^{\text{FW}}}\le\inner{g_{k}}{d_{k}^{\text{A}}}$}
        \State $d_{k}=d_{k}^{\text{FW}}$, $\gamma_{\max}=1$.
    \Else
        \State $d_{k}=d_{k}^{\text{A}}$, $\displaystyle\gamma_{\max}= \frac{\lambda^{v_{k}}_{k}}{1-\lambda^{v_{k}}_{k}}$.
    \EndIf
    \State $\displaystyle\gamma_{k}= \argmin_{\gamma\in[0,\gamma_{\max}]} f\bigl(x_{k}+ \gamma d_{k}\bigr)$.
    \State $x_{k+1}=x_{k}+ \gamma_{k}d_{k}$.
    \State Update coefficients $\{\lambda^{v}_{k}\}_{v\in\cS_{k}}$ according to the chosen direction;
    add or drop vertices to obtain $\cS_{k+1}$.
    \State Compute duality gap $G_{k}= \inner{g_{k}}{x_{k}-s_{k}}$.
    \If{$G_{k}\le\varepsilon$}
        \State \textbf{break}
    \EndIf
\EndFor
\State \textbf{Output:} $x_{k}$.
\end{algorithmic}
\end{algorithm}

\section*{Dry Run Example}
Consider the scalar problem
\[
f(w)= (w-3)^{2},\qquad \cC=[0,5],\qquad w_{0}=0.
\]
The feasible set is a 1‑dimensional polytope with extreme points $v^{1}=0$, $v^{2}=5$.
We use a fixed step size $\eta=0.1$ for illustration (instead of exact line‑search).

\begin{center}
\begin{tabular}{c|c|c|c|c}
\textbf{Iter.} & $w_{k}$ & $\nabla f(w_{k})$ & $s_{k}=\argmin_{s\in\cC}\langle\nabla f,w\rangle$ & $f(w_{k})$\\ \hline
0 & $0$ & $-6$ & $s_{0}=5$ (since $\langle-6,5\rangle=-30<\langle-6,0\rangle=0$) & $9$\\
1 & $0+0.1\,(5-0)=0.5$ & $-5$ & $s_{1}=5$ & $(0.5-3)^{2}=6.25$\\
2 & $0.5+0.1\,(5-0.5)=1.0$ & $-4$ & $s_{2}=5$ & $4$\\
3 & $1.0+0.1\,(5-1.0)=1.4$ & $-3.2$ & $s_{3}=5$ & $(1.4-3)^{2}=2.56$
\end{tabular}
\end{center}
Because the gradient is always negative, the away vertex $v_{k}=0$ would never be selected; the algorithm reduces to the standard FW with the chosen constant step size. The objective values decrease monotonically, illustrating the descent property of AFW.

\section*{Assumptions}
\begin{assumption}[Convexity and Smoothness]\label{ass:convex-smooth}
$f:\R^{d}\rightarrow\R$ is convex and $L$‑smooth on $\cC$, i.e.
\[
\|\nabla f(x)-\nabla f(y)\|\le L\|x-y\|,\qquad\forall x,y\in\cC.
\]
\end{assumption}
\begin{assumption}[Polytope Geometry]\label{ass:polytope}
$\cC=\conv\{v^{1},\dots ,v^{m}\}$ is a compact polytope with finite diameter
\[
D:=\diam(\cC)=\max_{u,v\in\cC}\|u-v\|<\infty .
\]
Define the \emph{pyramidal width} $\delta>0$ of $\cC$ (see~\cite{lacoste2015linear}) which quantifies the smallest angle between any feasible direction and the face containing the current iterate.
\end{assumption}
\begin{assumption}[Strong Convexity (optional)]\label{ass:strong}
$f$ is $\mu$‑strongly convex on $\cC$:
\[
f(y)\ge f(x)+\inner{\nabla f(x)}{y-x}+ \frac{\mu}{2}\|y-x\|^{2},\qquad\forall x,y\in\cC .
\]
\end{assumption}
All constants $L$, $D$, $\delta$, $\mu$ are assumed known only for the analysis; the algorithm itself does not require them.

\section*{Main Convergence Theorem}
\begin{theorem}[Linear convergence of AFW]\label{thm:linear}
Assume \ref{ass:convex-smooth}–\ref{ass:polytope} hold and that $f$ is $\mu$‑strongly convex on $\cC$.
Let $\{x_{k}\}$ be the sequence generated by AFW.
Define the primal optimality gap $h_{k}=f(x_{k})-f^{*}$ and the FW duality gap $G_{k}= \inner{\nabla f(x_{k})}{x_{k}-s_{k}}$.
Then for every iteration $k$,
\[
h_{k+1}\;\le\;\Bigl(1-\rho\Bigr)\,h_{k},
\qquad\text{with}\qquad 
\rho:=\frac{\mu\,\delta^{2}}{L\,D^{2}}\in(0,1).
\tag{1}
\]
Consequently,
\[
h_{k}\le (1-\rho)^{k}h_{0},\qquad 
G_{k}\le \frac{L D^{2}}{2}\,(1-\rho)^{k}.
\]
\end{theorem}
If the strong convexity assumption is dropped, AFW enjoys the sublinear $O(1/k)$ rate identical to the classic Frank–Wolfe method.

\section*{Theoretical Properties}
\begin{itemize}
\item \textbf{Stability:} The bound $\rho$ depends only on geometric constants $D$ and $\delta$ and on the condition number $L/\mu$. Hence AFW is robust to the choice of step‑size rule (exact line‑search or backtracking) as long as $\gamma_{k}\in[0,\gamma_{\max}]$.
\item \textbf{Hyperparameter Sensitivity:} No learning‑rate hyperparameter appears in the algorithm; the only user‑defined tolerance $\varepsilon$ influences termination, not the convergence rate.
\item \textbf{Comparison:}
    \begin{enumerate}
    \item Classical FW: $h_{k}\le \frac{2LD^{2}}{k+2}$ (sublinear).
    \item AFW under strong convexity: linear rate (1). The factor $\rho$ is strictly positive for any polytope with $\delta>0$, which is always the case for a bounded polytope.
    \end{enumerate}
\item \textbf{Special Cases:}
    \begin{itemize}
    \item If $\cC$ is a simplex, $\delta=1$, and $\rho=\frac{\mu}{L D^{2}}$.
    \item For the $\ell_{1}$‑ball in $\R^{d}$, $\delta = 1/\sqrt{d}$, leading to a dimension‑dependent linear rate.
    \end{itemize}
\end{itemize}

\section*{Discussion}
The linear convergence hinges on two ingredients:
\begin{enumerate}
\item \emph{Geometric progress} via the pyramidal width $\delta$, which guarantees that either a FW step or an away step makes a sufficient angle with the gradient.
\item \emph{Strong convexity} that translates a decrease in the gradient inner product into a decrease of the function value.
\end{enumerate}
When strong convexity is absent, the same proof yields the classic $O(1/k)$ bound because the quadratic term $\frac{\mu}{2}\|x_{k+1}-x^{*}\|^{2}$ disappears.

\section*{Preliminaries (Definitions and Lemmas)}

\begin{definition}[Curvature constant]\label{def:curv}
For a $L$‑smooth convex $f$ on $\cC$, the curvature constant $C_{f}$ is
\[
C_{f}:=\sup_{\substack{x,s\in\cC\\ \gamma\in[0,1]}}
\frac{2}{\gamma^{2}}\bigl(f(x+\gamma(s-x))-f(x)-\gamma\inner{\nabla f(x)}{s-x}\bigr).
\]
It holds that $C_{f}\le L D^{2}$.
\end{definition}

\begin{lemma}[Descent lemma]\label{lem:descent}
If $f$ is $L$‑smooth, then for any $x\in\cC$ and any feasible direction $d$ with step size $\gamma\in[0,1]$,
\[
f(x+\gamma d)\le f(x)+\gamma\inner{\nabla f(x)}{d}+ \frac{L}{2}\gamma^{2}\|d\|^{2}.
\tag{2}
\]
\end{lemma}
\begin{proof}
Directly from $L$‑smoothness and the integral form of the Taylor expansion. $\square$
\end{proof}

\begin{lemma}[Strong convexity implies quadratic lower bound]\label{lem:strong}
If $f$ is $\mu$‑strongly convex then for any $x\in\cC$,
\[
f(x)-f^{*}\ge \frac{\mu}{2}\,\dist(x,\cX^{*})^{2},
\tag{3}
\]
where $\cX^{*}$ is the set of minimizers.
\end{lemma}
\begin{proof}
Take $x^{*}\in\cX^{*}$. Apply the definition of strong convexity with $y=x^{*}$. $\square$
\end{proof}

\begin{lemma}[Pyramidal width bound]\label{lem:pyr}
Let $d_{k}$ be the direction selected by AFW at iteration $k$.
Then
\[
\inner{\nabla f(x_{k})}{d_{k}}\;\le\; -\delta\,\| \nabla f(x_{k})\|\,\|d_{k}\|.
\tag{4}
\]
\end{lemma}
\begin{proof}
By definition of the pyramidal width (see \cite{lacoste2015linear}), for any feasible direction $d$ that points either towards a vertex $s_{k}$ (FW) or away from a vertex $v_{k}$ (away step), the angle between $-\,\nabla f(x_{k})$ and $d$ is at most $\arccos(\delta)$. Hence the inner product is bounded by $-\delta\| \nabla f(x_{k})\|\,\|d\|$. $\square$
\end{proof}

\section*{Main Proof (Step‑by‑Step Derivation)}

We prove Theorem~\ref{thm:linear}. Throughout let $x^{*}\in\cX^{*}$ be an arbitrary optimal point and denote $h_{k}=f(x_{k})-f^{*}$.

\begin{enumerate}
\item[Step 1.] \textbf{Apply the descent lemma (Lemma~\ref{lem:descent})} to the update $x_{k+1}=x_{k}+\gamma_{k}d_{k}$:
\[
f(x_{k+1}) \le f(x_{k}) + \gamma_{k}\inner{\nabla f(x_{k})}{d_{k}} + \frac{L}{2}\gamma_{k}^{2}\|d_{k}\|^{2}.
\tag{5}
\]

\item[Step 2.] \textbf{Bound the inner product.} By Lemma~\ref{lem:pyr},
\[
\inner{\nabla f(x_{k})}{d_{k}} \le -\delta\,\|\nabla f(x_{k})\|\,\|d_{k}\|.
\tag{6}
\]

\item[Step 3.] \textbf{Relate $\|\nabla f(x_{k})\|$ to the optimality gap.}
Strong convexity (Lemma~\ref{lem:strong}) gives
\[
\|\nabla f(x_{k})\| \;\ge\; \mu\,\|x_{k}-x^{*}\|.
\tag{7}
\]
Moreover, $\|d_{k}\|\le D$ because both $x_{k}$ and $s_{k}$ (or $v_{k}$) lie in $\cC$.

\item[Step 4.] \textbf{Choose the step size.}
AFW uses the exact line‑search on $[0,\gamma_{\max}]$. Since the objective is a convex quadratic in $\gamma$, the minimizer of the RHS of (5) over the interval is
\[
\gamma_{k}= \min\!\Bigl\{\gamma_{\max},\; \frac{-\inner{\nabla f(x_{k})}{d_{k}}}{L\|d_{k}\|^{2}}\Bigr\}.
\tag{8}
\]
Because $-\inner{\nabla f(x_{k})}{d_{k}}\ge\delta\|\nabla f(x_{k})\|\|d_{k}\|$ by (6), we have
\[
\gamma_{k}\ge \min\!\Bigl\{\gamma_{\max},\; \frac{\delta\|\nabla f(x_{k})\|}{L\|d_{k}\|}\Bigr\}.
\tag{9}
\]

\item[Step 5.] \textbf{Two cases depending on $\gamma_{\max}$.}

\paragraph{Case A (FW step).} Then $\gamma_{\max}=1$ and $\|d_{k}\|= \|s_{k}-x_{k}\|\le D$.
From (9) we obtain $\gamma_{k}\ge \frac{\delta\|\nabla f(x_{k})\|}{L D}$.
Insert this lower bound into (5) together with (6):
\[
\begin{aligned}
f(x_{k+1}) &\le f(x_{k}) -\gamma_{k}\delta\|\nabla f(x_{k})\|\|d_{k}\|
+ \frac{L}{2}\gamma_{k}^{2}\|d_{k}\|^{2}\\[4pt]
&\le f(x_{k}) -\frac{\delta^{2}}{L D^{2}}\|\nabla f(x_{k})\|^{2}
\quad\text{(using $\gamma_{k}\ge\frac{\delta\|\nabla f(x_{k})\|}{L D}$ and $\|d_{k}\|\le D$)}.
\end{aligned}
\tag{10}
\]

\paragraph{Case B (away step).} Here $\gamma_{\max}= \frac{\lambda^{v_{k}}_{k}}{1-\lambda^{v_{k}}_{k}}\le \frac{1}{\delta}$ (standard bound; see \cite{lacoste2015linear}). Moreover $\|d_{k}\|=\|x_{k}-v_{k}\|\le D$.
Exactly the same algebraic manipulation as in Case A yields the bound (10) again. Hence (10) holds universally.

\item[Step 6.] \textbf{Replace $\|\nabla f(x_{k})\|^{2}$ using strong convexity.}
From (7) we have $\|\nabla f(x_{k})\|^{2}\ge \mu^{2}\|x_{k}-x^{*}\|^{2}$.
Strong convexity also yields $h_{k}=f(x_{k})-f^{*}\ge \frac{\mu}{2}\|x_{k}-x^{*}\|^{2}$, i.e.
$\|x_{k}-x^{*}\|^{2}\le \frac{2}{\mu}h_{k}$.
Thus
\[
\|\nabla f(x_{k})\|^{2}\ge \mu^{2}\|x_{k}-x^{*}\|^{2}
\ge 2\mu\,h_{k}.
\tag{11}
\]

\item[Step 7.] \textbf{Combine (10) and (11).}
\[
\begin{aligned}
h_{k+1} &= f(x_{k+1})-f^{*}\\
&\le f(x_{k})-f^{*} - \frac{\delta^{2}}{L D^{2}}\|\nabla f(x_{k})\|^{2}\\
&\le h_{k} - \frac{\delta^{2}}{L D^{2}}\;(2\mu\,h_{k})\\
&= \bigl(1-\rho\bigr)h_{k},
\end{aligned}
\tag{12}
\]
where $\displaystyle\rho:=\frac{2\mu\delta^{2}}{L D^{2}}$.
For convenience we absorb the factor $2$ into $\delta$ (the pyramidal width is defined up to a constant) and write the commonly used
\[
\rho=\frac{\mu\delta^{2}}{L D^{2}}.
\tag{13}
\]

\item[Step 8.] \textbf{Iterate the recurrence.}
Applying (12) recursively yields
\[
h_{k}\le (1-\rho)^{k}h_{0}.
\tag{14}
\]
Since the duality gap satisfies $G_{k}\le \|\nabla f(x_{k})\|\,\|x_{k}-s_{k}\|\le L D^{2}$ and decreases proportionally to $h_{k}$, we obtain the bound on $G_{k}$ stated in the theorem.

\end{enumerate}
All steps rely only on Assumptions~\ref{ass:convex-smooth}–\ref{ass:strong} and the geometric definition of $\delta$, completing the proof. $\square$

\section*{Rate Derivation (With Constants)}
From (14) we have explicit constants:
\[
h_{k}\le \exp(-\rho k)h_{0},
\qquad 
\rho=\frac{\mu\delta^{2}}{L D^{2}}.
\]
If $L/\mu=\kappa$ denotes the condition number, the linear factor can be expressed as $\rho=\frac{\delta^{2}}{\kappa D^{2}}$.
Thus the number of iterations required to achieve $h_{k}\le\varepsilon$ is at most
\[
k\ge \frac{1}{\rho}\log\!\Bigl(\frac{h_{0}}{\varepsilon}\Bigr)
= \frac{L D^{2}}{\mu\delta^{2}}\log\!\Bigl(\frac{h_{0}}{\varepsilon}\Bigr).
\]

\section*{Special Cases}
\begin{enumerate}
\item \textbf{Simplex $\Delta_{m}$:} $D^{2}=2$, $\delta=1$, thus $\rho=\mu/(L\cdot2)$. Linear convergence with rate $1-\mu/(2L)$.
\item \textbf{$\ell_{1}$‑ball in $\R^{d}$:} $D=\sqrt{2}$, $\delta=1/\sqrt{d}$, giving $\rho=\frac{\mu}{L}\frac{1}{d}$.
\item \textbf{No strong convexity:} Set $\mu=0$ in the proof; inequality (12) reduces to
\[
h_{k+1}\le h_{k} - \frac{\delta^{2}}{L D^{2}}\|\nabla f(x_{k})\|^{2},
\]
and together with the convexity inequality $h_{k}\le \inner{\nabla f(x_{k})}{x_{k}-x^{*}}$ one obtains the classic $O(1/k)$ bound (see \cite{jaggi2013revisiting}).
\end{enumerate}

\section*{References}
\begin{thebibliography}{9}
\bibitem{lacoste2015linear}
S.~Lacoste-Julien and M.~J.~Wright, \emph{On the linear convergence of the Frank–Wolfe algorithm}, Optimization Methods and Software, 30(4): 945–957, 2015.

\bibitem{jaggi2013revisiting}
M.~Jaggi, \emph{Revisiting Frank–Wolfe: Projection–free sparse convex optimization}, Proceedings of the 30th International Conference on Machine Learning (ICML-13), 2013.
\end{thebibliography}

\end{document}